% Construcción de imágenes con Podman
\subsection{Construcción de imágenes con Podman}\label{sec:4_3_construccion_imagenes}
Una imagen es la plantilla a partir de la cual Podman crea un contenedor; se describe en un archivo \texttt{Dockerfile}. La plataforma usa tres Dockerfile. Las imágenes \textbf{no} se construyen en el equipo de desarrollo: se construyen en cada VPS con sus scripts (sección \ref{sec:4_3_4_build}), a partir de los artefactos que se transfieren en la sección \ref{sec:4_4_transferencia_vps}.

\subsubsection{Dockerfile de los microservicios (eclipse-temurin:21-jre-jammy)}
\texttt{Backend/Dockerfile} es común a todos los microservicios y al modelo de IA:
\begin{lstlisting}
FROM docker.io/library/eclipse-temurin:21-jre-jammy
WORKDIR /app
ARG JAR_FILE
ARG SERVICE_PORT=8080
COPY ${JAR_FILE} app.jar
EXPOSE ${SERVICE_PORT}
ENTRYPOINT ["java", "-jar", "app.jar"]
\end{lstlisting}
\begin{itemize}
  \item \texttt{FROM}: parte de una imagen oficial que solo trae el entorno de ejecución de Java 21 (JRE), más ligera que un JDK completo.
  \item \texttt{ARG JAR\_FILE}: el JAR se recibe como parámetro al construir, por eso un solo Dockerfile sirve para los once servicios.
  \item \texttt{COPY}: copia el JAR dentro de la imagen como \texttt{/app/app.jar}.
  \item \texttt{ENTRYPOINT}: al arrancar el contenedor ejecuta \texttt{java -jar app.jar}.
\end{itemize}

\subsubsection{Dockerfile de admin-service (cliente pg\_dump 16)}
\texttt{admin-service} genera y restaura los respaldos del Panel de la Defensora con \texttt{pg\_dump} y \texttt{psql}, así que su imagen necesita además el cliente de PostgreSQL. Ubuntu Jammy (la base de \texttt{eclipse-temurin}) trae la versión 14 y \texttt{pg\_dump} 14 se niega a respaldar un servidor PostgreSQL 16, por lo que \texttt{Backend/admin-service/Dockerfile} instala la versión 16 desde el repositorio oficial de PostgreSQL:
\begin{lstlisting}
FROM docker.io/library/eclipse-temurin:21-jre-jammy
WORKDIR /app
RUN apt-get update \
 && apt-get install -y --no-install-recommends curl ca-certificates \
 && (se agrega el repositorio apt.postgresql.org con su llave) \
 && apt-get install -y --no-install-recommends postgresql-client-16 \
 && rm -rf /var/lib/apt/lists/* && pg_dump --version
ARG JAR_FILE
COPY ${JAR_FILE} app.jar
RUN mkdir -p /app/respaldos
ENTRYPOINT ["java", "-jar", "app.jar"]
\end{lstlisting}
La carpeta \texttt{/app/respaldos} se monta como volumen (\nolinkurl{/apps/utiles/respaldos}) para que los respaldos sobrevivan a la reconstrucción del contenedor.

\subsubsection{Dockerfile de los frontends (nginx:alpine)}
Los frontends se sirven con un nginx mínimo que solo entrega los archivos estáticos:
\begin{lstlisting}
FROM docker.io/library/nginx:alpine
COPY nginx.conf /etc/nginx/conf.d/default.conf
COPY dist/browser /usr/share/nginx/html
EXPOSE 80
\end{lstlisting}
La configuración interna de nginx de cada frontend contiene \texttt{try\_files \$uri \$uri/ /index.html}, para que al recargar una página interna de Angular (por ejemplo \texttt{/admin/personal}) se devuelva \texttt{index.html} y Angular resuelva la ruta. El portal (\texttt{Frontend}) usa la misma estructura con su archivo \texttt{config/static.conf} y escucha en el 8090 en lugar del 80.

\captura{comp-09-dockerfiles}{Contenido de los Dockerfile de microservicios y frontends}

\subsubsection{Construcción de las imágenes en el servidor con los scripts podman-compose}\label{sec:4_3_4_build}
Cada script construye la imagen con \texttt{podman build} y crea el contenedor con \texttt{podman run}. Por ejemplo, para \texttt{quejas-service}, \texttt{podman-compose.sh} ejecuta:
\begin{lstlisting}
podman build -f Dockerfile --build-arg JAR_FILE=artifact/quejas-service.jar \
             -t defensoria-quejas-service .
\end{lstlisting}
\begin{table}[H]
\centering
\caption{Imágenes que construye cada script}
\small
\begin{tabular}{p{2.2cm}p{4.6cm}p{7.6cm}}
\toprule
\textbf{VPS} & \textbf{Script} & \textbf{Imágenes} \\
\midrule
Backend & \texttt{podman-compose.sh} & \texttt{defensoria-<servicio>} (una por microservicio) \\
Backend & \texttt{podman-compose-modelo.sh} & \texttt{defensoria-antecedentes-service} \\
Frontend & \texttt{podman-compose-front*.sh} & \texttt{defensoria-front-img}, \texttt{defensoria-admin-img}, \texttt{defensoria-revision-img}, \texttt{defensoria-primer-contacto-img}, \texttt{defensoria-subdefensoria-img} \\
Base de datos & \texttt{podman-compose-bd.sh} & No construye: usa directamente \texttt{postgres:16-alpine} \\
\bottomrule
\end{tabular}
\end{table}
La ejecución de cada script se describe en los capítulos 5, 6 y 7.
